% context: 9th-grade physics exit note, third document of the year
% topic: rounding to 3 s.f., scientific notation both directions, unit conversion, mi/h to ft/s
% structure: front page is 5-10 minutes of short work; back page is one long challenge
\documentclass[12pt, twoside]{article}
\input{../preamble}
\usepackage{float}

\title{Physics}
\author{Chris Huson}
\date{September 2026}

\fancyhead[LE]{\thepage}
\fancyhead[RO]{\thepage \\ Nickname: \hspace{2.25cm} \,\\ \,}
\fancyhead[LO]{La Scuola d'Italia / Huson / Physics \\* 18 September 2026}

\begin{document}

\subsubsection*{1.3 Exit Note: Unit Conversions}

\textbf{Reference (provided).} \quad
1 mi = 5280 ft \quad 1 m $\approx$ 3.28 ft \quad 1 ft = 12 in \\
1 km = 1000 m \quad 1 h = 3600 s \quad 1 min = 60 s

\vspace{0.35cm}

\begin{enumerate}[itemsep=0.7cm]

\item Round each value to three significant figures. Copy the calculator display
  followed by three dots, then round: $\sqrt{2} = 1.4142135\ldots \approx 1.41$
  \begin{multicols}{2}
    \begin{enumerate}[itemsep=1.5cm]
      \item $\sqrt{3}= 1.7320508076...$
      \item $\dfrac{100~\text{m}}{9.58~\text{s}}$ \quad (Usain Bolt's average speed
        in the 100 m world record)
    \end{enumerate}
  \end{multicols} \vspace{0.5cm}

\item Scientific notation.
  \begin{enumerate}[itemsep=2.2cm]
    \item One year is 31{,}556{,}952 seconds. Write this in scientific notation,
      rounded to three significant figures.
    \item A hydrogen atom is $1.06 \times 10^{-10}$ meters across. Write this as an
      ordinary decimal number.
  \end{enumerate} \vspace{2.0cm}

\item The Duomo di Milano is 108 meters tall. What is its height in feet? Write the
  conversion as a fraction and show the units cancelling. \vspace{2cm}

\item A Category 1 hurricane begins at a wind speed of 74 miles per hour. What is
  this speed in feet per second? \vspace{3.4cm}

\newpage

\item \textbf{Challenge: how fast is an elevator?}

  The elevator industry does not quote speeds in meters per second. Each
  manufacturer uses the unit customary where the machine was built:

  \begin{center}
  \begin{tabular}{l l}
    \textbf{Building} & \textbf{Advertised speed} \\[0.12cm]
    \hline \\[-0.28cm]
    Shanghai Tower & 20.5 meters per second \\[0.1cm]
    Taipei 101 & 1010 meters per minute \\[0.1cm]
    One World Trade Center & 2000 feet per minute \\[0.1cm]
    Empire State Building (express) & 1200 feet per minute \\
  \end{tabular}
  \end{center}

  \vspace{0.15cm}

  \begin{enumerate}[itemsep=0.35cm]
    \item Convert all four speeds to meters per second, rounded to three
      significant figures. Write each one as a single chain of fractions and show
      the units cancelling. \vspace{4.0cm}

    \item List the four elevators from slowest to fastest. Is that ranking
      supported by the precision of the data? Explain. \vspace{1.7cm}

    \item Our classroom is 18 meters above the building entrance. If each of these
      elevators were installed in our building and ran at its advertised speed, how
      long would the ride up take? Compare with the 75 seconds the stairs took us. \vspace{2.1cm}

    \item A real elevator has to speed up at the start of the ride and slow down at
      the end. Explain why your answers to part (c) are all too small --- and why
      the error is worst for the Shanghai Tower.
  \end{enumerate}

\end{enumerate}
\end{document}
